\honest{Humans Are Not Automatically Strategic}{Anna Salamon}{2010}

I am replying to a post by Lionhearted, who asks why a man who wants to be the best comedian ever would study by watching reruns of \textsc{Garfield and Friends}. The comedian is my only worked example.

My answer starts with a child. An eight-year-old fails a calculus test because most possible answers are wrong and nothing guides the child to the right ones. Failure is the default, and needs no special explanation such as a ``fear of success''; people, and rocks, fail calculus tests by default.

I apply this to adults. Most ways of pursuing a goal are ineffective, and no evolutionary or cultural force has pushed us toward the few that work. I call this a guess, and give no evidence for the part about evolution and culture. The other examples go in a footnote: people train for a career without first comparing salaries, type with two fingers for years, never check which of their hobbies they enjoy, fear illness and accidents without looking up the relative risks, and keep one way of studying or writing without trying others.

I mention only in passing that the adult, unlike the child, may be unsure whether the effort is worth it. \nb{Looking things up takes time. Stopping at good enough is often sensible; Herbert Simon called it satisficing. A two-finger typist near retirement may be right not to learn.}

I grant that we have goals in three weak senses: we tell stories about them, we do things that fit the role, and we feel glad or disappointed at the result. Then I list eight things we do not do automatically. We do not ask what we want, ask how we would know we had got it, feel curious about how to get it, gather that information, test different methods, focus on what works, make sure the goal ``is really our goal,'' free of fears and of doubt whether it is worth the effort, with decisions thought through in advance so they do not sap our energy, or use our surroundings and social settings to keep us motivated through frustration and temptations from hyperbolic discounting. Instead ``we mostly just do things.'' We act from habit or impulse, or pick an action that ``feels associated'' with the goal.

\nb{A simpler explanation of the comedian is that he likes cartoons more than he wants the hard work, and that ``becoming a comedian'' is partly a story he tells. Two items in the post come near this: goals as stories we tell, and making sure the goal ``is really our goal.''} On my account we have goals only in ``limited senses''; yet we do not optimize ``for any other goals'' either, so I treat the gap between what people say they want and what they do as a lack of skill.

Why do we not do these things? Because ``humans are only just on the cusp of general intelligence.'' I give nothing to support this. \nb{The rest of the argument does not need it.} Perhaps 5\% of people can even understand these methods once they are pointed out. I give no source for the figure. Our verbal systems reason abstractly far better than the motivational systems that pull our behavior. Knowing that exercise is good for you does not make you want to exercise. I know that a glass floor high up in a building is safe, and I still have to walk across it many times before the rest of my brain agrees.

Training those systems, say by picturing ice cream as yucky, is not automatic either, so it is no surprise that most of our goal-seeking is ineffective. I am keen to train; I know people far more strategic than I am. There ``seem to be clear avenues'' for becoming more strategic, and I do not name one beyond my own list and an exercise of imagining ice cream as disgusting. Half a sentence adds that having goals in this fuller sense is part of what ``rational'' should mean, and has not been taught in much detail on LW. Last, I ask readers how they have trained themselves.

\nb{Two years later Salamon co-founded the Center for Applied Rationality, which teaches this kind of training in paid workshops.}
